% =====================================================================
\section{Expected information gain in search and localize}
\label{sm:eig}
% =====================================================================

Sec.~\ref{sec:eig} of the main text asks two questions of the variance form of the expected information gain,
Eq.~\eqref{eq:eig-var}: (A) when it agrees with the exact gain, and (B) when maximizing it greedily puts the
measurements where the van Trees bound needs them. In refine the answer to both is yes. This section gives the
argument for the two earlier regimes, where both fail.

\paragraph{Search: a two-valued signal.} Before the dip is found the posterior of $\theta$ is spread over the window,
and the predicted signal at $x$ is two-valued: with probability $p$, the posterior mass that would put the dip at $x$,
it is about $1-C$, otherwise 1. Then $\sigma^2_{\rm pred}=p(1-p)C^2$, while the true mutual information is at most the
binary entropy $H_b(p)\le\ln2$. For (A), the variance form keeps growing with SNR where the true gain saturates, so it
overvalues search candidates relative to refine candidates; both peak at $p=\tfrac12$, so they rank search
candidates alike. For (B), the chain-rule argument of Eq.~\eqref{eq:eig-sum} fails, since where one measures changes
what later measurements can gain. Here greedy EIG is only myopically optimal, which is where the strategies of
Sec.~\ref{sec:locators} of the main text differ.

\paragraph{Localize: the variance form chases the tail.} The same two-valued picture holds in the localize stage,
with a different ranking. Once a peak has formed at the dip, a small residual mass $p$ far from it gives a predictive
variance $p(1-p)C^2$ at its location, large at high SNR even for small $p$, while the information it holds is at most
$H_b(p)$, small for small $p$. The variance form therefore keeps spending measurements to remove such distant mass.
This reduces the posterior variance, which the tail dominates (Sec.~\ref{sm:gaussian}) and which is the quantity
Eq.~\eqref{eq:vantrees} constrains, rather than the posterior entropy: in the localize stage the variance rule acts
as a design for the MSE, while a rule built on the information gain would refine the peak sooner and accept the
tail. Whether this explains why localizing costs SBED far more measurements than the bound is a hypothesis we have
not tested.
